\documentclass[10pt,twocolumn]{article}

\usepackage[margin=0.75in,top=0.70in,bottom=0.70in]{geometry}
\usepackage{times}
\usepackage{amsmath,amssymb}
\usepackage{booktabs}
\usepackage{graphicx}
\usepackage{algorithmic}
\usepackage{algorithm}
\usepackage{hyperref}
\usepackage{microtype}
\usepackage{xcolor}
\usepackage{natbib}
\usepackage{enumitem}
\setlist{nosep,leftmargin=*}

% Tight spacing for conference style
\usepackage{titlesec}
\titlespacing*{\section}{0pt}{1.0ex plus 0.2ex minus 0.2ex}{0.4ex plus 0.1ex}
\setlength{\parskip}{0pt}
\setlength{\parindent}{1em}
\setlength{\textfloatsep}{8pt plus 2pt minus 2pt}
\setlength{\floatsep}{8pt plus 2pt minus 2pt}
\setlength{\intextsep}{6pt plus 2pt minus 2pt}
\setlength{\abovecaptionskip}{4pt}
\setlength{\belowcaptionskip}{2pt}

\title{\Large\bfseries Bidirectional Message Passing for Coupled Cardiorenal Hemodynamic Simulation}
\author{Pranav Dorbala\\
\normalsize\itshape Johns Hopkins University, Baltimore, MD}
\date{}

\begin{document}
\maketitle

\begin{abstract}
\noindent We present a coupled cardiorenal simulator that models bidirectional physiological feedback between cardiac and renal subsystems via structured message passing. A time-varying elastance heart model with an exponential end-diastolic pressure-volume relationship (EDPVR) exchanges hemodynamic signals with a mechanistic renal model (Hallow et al., 2017) implementing tubuloglomerular feedback (TGF), the renin-angiotensin-aldosterone system (RAAS), and tubular sodium handling. The coupling architecture transmits three cardiac-to-renal messages (mean arterial pressure, cardiac output, central venous pressure) and two renal-to-cardiac messages (blood volume, systemic vascular resistance) at each discrete time step, governed by a coupling intensity parameter $\alpha$. We introduce a counterfactual decomposition that attributes each organ's deterioration rate to three mechanistic drivers: its own disease state, the other organ's current state, and the other organ's rate of decline. The framework reproduces four cardiorenal syndrome phenotypes and quantifies the amplification of deterioration rates attributable to cross-organ coupling.
\end{abstract}

\section{Problem Formulation}
\label{sec:formulation}

Cardiorenal syndrome (CRS) describes disorders in which acute or chronic dysfunction of the heart or kidneys induces dysfunction in the other organ \citep{ronco2008}. We formulate this as a \emph{coupled dynamical system} over discrete time steps $t = 1, \dots, T$. Let $\mathbf{h}_t$ denote the cardiac state (ventricular volumes, pressures, ejection fraction) and $\mathbf{r}_t$ the renal state (glomerular filtration rate, blood volume, sodium balance). Each subsystem advances using its own dynamics \emph{and} messages from the other:
\begin{align}
\mathbf{h}_{t+1} &= f_H\!\bigl(\mathbf{h}_t,\; \mathbf{m}_{t}^{K \to H},\; \theta_t^H\bigr) \label{eq:heart}\\
\mathbf{r}_{t+1} &= f_K\!\bigl(\mathbf{r}_t,\; \mathbf{m}_{t}^{H \to K},\; \theta_t^K\bigr) \label{eq:kidney}
\end{align}
where $\theta_t^H$ and $\theta_t^K$ are time-varying disease parameters (myocardial stiffness and ultrafiltration coefficient), and the message vectors are:
\begin{align}
\mathbf{m}^{H \to K}_t &= \bigl[\text{MAP}_t,\; \text{CO}_t,\; \text{CVP}_t\bigr]^\top \label{eq:h2k}\\
\mathbf{m}^{K \to H}_t &= \bigl[V_{\text{blood},t},\; \text{SVR}_t\bigr]^\top \label{eq:k2h}
\end{align}

A scalar coupling intensity $\alpha \in \{0, 1, 2\}$ modulates every message:
\begin{equation}
\hat{m}_i = m_i^{(0)} + \alpha \cdot \bigl(m_i - m_i^{(0)}\bigr)
\label{eq:alpha}
\end{equation}
where $m_i^{(0)}$ is the healthy baseline. Setting $\alpha = 0$ decouples the organs; $\alpha = 1$ recovers normal physiology; $\alpha > 1$ amplifies the feedback.


\section{Cardiac Subsystem}
\label{sec:heart}

The left ventricle is modeled as a time-varying elastance chamber with an exponential EDPVR capturing HFpEF diastolic dysfunction. End-diastolic volume (EDV) is determined by inverting the EDPVR at the current left atrial filling pressure $P_{\text{LA}}$:
\begin{equation}
P_{\text{LA}} = \gamma \cdot A_{\text{ed}}\bigl(e^{\,\beta\,(V_{\text{ED}} - V_0)} - 1\bigr)
\label{eq:edpvr}
\end{equation}
where $\gamma \geq 1$ is the \textbf{stiffness scale} (the primary HFpEF parameter), $A_{\text{ed}} = 0.5$\,mmHg, $\beta = 0.028$\,mL$^{-1}$, and $V_0 = 10$\,mL. Filling pressure derives from blood volume via simplified venous compliance:
\begin{equation}
P_{\text{LA}} = 10 + \frac{V_{\text{blood}} - 5000}{200} \quad \text{(mmHg)}
\label{eq:pla}
\end{equation}
End-systolic volume is set by maximum elastance $E_{\max} = 2.3$\,mmHg/mL against effective systemic resistance $R_{\text{eff}} = R_{\text{sys}} \cdot R_{\text{ratio}} \cdot \xi$, where $R_{\text{ratio}}$ is the renal-driven SVR message and $\xi$ captures arterial stiffness. Stroke volume, ejection fraction, cardiac output ($\text{CO} = \text{SV} \times \text{HR}$), and MAP follow.

\textbf{Key behavior.} Because the EDPVR is exponential, stiffness interacts \emph{multiplicatively} with blood volume: a stiff ventricle ($\gamma > 1$) at elevated $V_{\text{blood}}$ produces disproportionately elevated LVEDP---the hallmark of HFpEF decompensation under volume overload.


\section{Renal Subsystem}
\label{sec:kidney}

The kidney follows Hallow et al.\ \citeyearpar{hallow2017}, a mechanistic model of glomerular hemodynamics and tubular transport.

\textbf{RAAS.} The renin-angiotensin-aldosterone system responds to MAP deviation from a homeostatic setpoint:
\begin{equation}
\text{RAAS} = \text{clip}\bigl(1 - g_R \cdot 0.005 \cdot \Delta\text{MAP},\; 0.5,\; 2.0\bigr)
\label{eq:raas}
\end{equation}
This modulates efferent arteriolar resistance $R_{\text{EA}} = R_{\text{EA}}^{(0)} \cdot \text{RAAS}$ and collecting duct sodium reabsorption $\eta_{\text{CD}} = \eta_{\text{CD}}^{(0)} \cdot \text{RAAS}$. The RAAS factor also feeds back as part of the SVR message to the heart: elevated RAAS drives systemic vasoconstriction.

\textbf{Tubuloglomerular feedback.} An iterative solver (12 fixed-point iterations) equilibrates afferent arteriolar resistance $R_{\text{AA}}$ based on macula densa sodium delivery:
\begin{equation}
R_{\text{AA}}^{(k+1)} \!=\! R_{\text{AA}}^{(0)} \!\Bigl(\!1 + g_{\text{TGF}} \!\cdot\! \frac{\text{MD}_{\text{Na}}^{(k)} - \text{MD}_{\text{Na}}^{*}}{\text{MD}_{\text{Na}}^{*}}\!\Bigr)
\label{eq:tgf}
\end{equation}
Glomerular capillary pressure, renal blood flow, and single-nephron GFR are simultaneously determined:
\begin{align}
\text{RBF} &= \frac{\text{MAP} - P_{\text{ven}}}{R_{\text{pre}} + R_{\text{AA}} + R_{\text{EA}}} \cdot 10^3 \label{eq:rbf}\\
P_{\text{gc}} &= \text{MAP} - \text{RBF}(R_{\text{pre}} \!+\! R_{\text{AA}}) / 10^3 \label{eq:pgc}\\
\text{SNGFR} &= K_f \cdot K_f^{\text{scale}} \cdot \text{NFP} \label{eq:sngfr}
\end{align}
where $\text{NFP} = P_{\text{gc}} - P_{\text{Bow}} - \bar{\pi}$ is net filtration pressure, $K_f^{\text{scale}} \leq 1$ is the CKD disease parameter (nephron loss), and total $\text{GFR} = 2N \cdot \text{SNGFR}$.

\textbf{Volume balance.} Filtered sodium traverses four tubular segments with fractional reabsorption, modulated by pressure natriuresis. The intake--excretion difference for sodium and water updates $V_{\text{blood}}$ over $\Delta t$.

\section{Coupling Architecture}
\label{sec:coupling}

Algorithm~\ref{alg:coupling} details the message-passing loop. At each step, the heart produces hemodynamics from stiffness $\gamma_t$ and the kidney's last-reported $V_{\text{blood}}$ and SVR. Three $\alpha$-scaled messages are delivered to the kidney, which advances its state and returns updated volume and resistance.

\begin{algorithm}[t]
\caption{Coupled Cardiorenal Simulation}\label{alg:coupling}
\begin{algorithmic}[1]
\small
\STATE \textbf{Input:} $\{\gamma_t\}$, $\{K_f^{\text{s}}_t\}$, $\alpha$, $\Delta t$, $T$
\STATE Init.\ $\mathbf{r}_0$; pre-equilibrate TGF setpoint
\STATE $V_b \!\leftarrow\! 5000$;\; $R_r \!\leftarrow\! 1.0$;\; store baselines step 1
\FOR{$t = 1$ \TO $T$}
    \STATE $\mathbf{h}_t \leftarrow f_H(\gamma_t,\, V_b,\, R_r)$ \hfill\COMMENT{Cardiac model}
    \STATE $\widehat{\text{MAP}},\, \widehat{\text{CO}},\, \widehat{\text{CVP}} \leftarrow$ scale by $\alpha$ \hfill\COMMENT{Eq.~\ref{eq:alpha}}
    \STATE $\mathbf{r}_t \leftarrow f_K(\mathbf{r}_{t\!-\!1},\, \widehat{\text{MAP}},\, \widehat{\text{CO}},\, \widehat{\text{CVP}},\, \Delta t)$
    \STATE $\text{SVR}_t \leftarrow \text{RAAS}_t \cdot (1 + 0.4\,(1 \!-\! K_f^{\text{s}}_t))$
    \STATE $R_r \leftarrow 1 + \alpha\,(\text{SVR}_t - 1)$;\enspace $V_b \leftarrow \mathbf{r}_t.V_b$
\ENDFOR
\end{algorithmic}
\end{algorithm}

The SVR message (line~8) is composed of two renal pathology signals: RAAS-mediated vasoconstriction (neurohormonal) and CKD-driven vascular stiffening ($0.4 \cdot (1 - K_f^{\text{scale}})$). This dual pathway captures both the acute hemodynamic and chronic structural effects of kidney disease on cardiac afterload.

\section{Counterfactual Decomposition}
\label{sec:decomposition}

To attribute each organ's deterioration to specific mechanistic drivers, we decompose the per-step change $\Delta y_t = y_t - y_{t-1}$ for a target variable $y$ into three additive components.

\textbf{Cardiac decomposition} ($y = \text{EF}$). Let $\text{EF}(\gamma, V, R)$ denote the heart model's ejection fraction at given stiffness, volume, and resistance. At step~$t$:
\begin{align}
\delta_H &= \text{EF}(\gamma_t, V_0, 1) - \text{EF}(\gamma_{t\!-\!1}, V_0, 1) \label{eq:dh}\\
\delta_{K_s} &= \text{EF}(\gamma_t, V_{t\!-\!1}, R_{t\!-\!1}) - \text{EF}_{t\!-\!1} - \delta_H \label{eq:dks}\\
\delta_{K_r} &= \Delta\text{EF}_t - \delta_H - \delta_{K_s} \label{eq:dkr}
\end{align}
where $\delta_H$ isolates the effect of stiffness progression at baseline renal conditions (\textbf{own disease}), $\delta_{K_s}$ captures amplification from the kidney's \emph{accumulated} pathological state---elevated $V_{\text{blood}}$ and SVR (\textbf{other organ's state}), and $\delta_{K_r}$ is the residual attributable to the kidney's \emph{change this step}---$\Delta V_{\text{blood}}$ and $\Delta$SVR (\textbf{other organ's rate of decline}).

\textbf{Renal decomposition} ($y = \text{GFR}$). A parallel uncoupled simulation ($\alpha = 0$) with identical disease schedules yields the counterfactual:
\begin{align}
\delta_K &= \Delta\text{GFR}_t^{\,\alpha=0} \quad \text{(own disease)} \label{eq:dk}\\
\delta_{H_{\text{tot}}} &= \Delta\text{GFR}_t - \delta_K \quad \text{(cardiac influence)} \label{eq:dhtot}
\end{align}
The cardiac influence $\delta_{H_{\text{tot}}}$ is split between heart \emph{state} and heart \emph{rate} by the ratio of accumulated hemodynamic deficit (MAP depression, CVP elevation relative to normals) to instantaneous hemodynamic change ($|\Delta\text{MAP}_t|$, $|\Delta\text{CVP}_t|$):
\begin{align}
w &= \frac{|\Delta\text{MAP}_t| + |\Delta\text{CVP}_t|}{D_{\text{acc}} + |\Delta\text{MAP}_t| + |\Delta\text{CVP}_t|} \nonumber\\[3pt]
\delta_{H_s} &= (1\!-\!w)\,\delta_{H_{\text{tot}}}, \qquad \delta_{H_r} = w\,\delta_{H_{\text{tot}}} \label{eq:dhsplit}
\end{align}
where $D_{\text{acc}} = |\overline{\text{MAP}} - \text{MAP}_t| + \max(\text{CVP}_t - 3, 0)$ is the accumulated hemodynamic deficit. This provides a per-step attribution of GFR decline to intrinsic nephron loss ($\delta_K$), the heart's current hemodynamic deficit ($\delta_{H_s}$), and the heart's ongoing deterioration ($\delta_{H_r}$).

\textbf{Rate amplification.} The aggregate amplification due to coupling is:
\begin{equation}
\mathcal{A}_y = \frac{\overline{|\Delta y^{\alpha}|} - \overline{|\Delta y^{0}|}}{\max\!\bigl(\overline{|\Delta y^{0}|},\; \tfrac{1}{2}\overline{|\Delta y^{\alpha}|},\; \epsilon\bigr)}
\label{eq:amplification}
\end{equation}
where bars denote time-averaged absolute per-step rates and $\epsilon$ prevents division by zero when the uncoupled trajectory is near-stationary.

\section{Scenarios}
\label{sec:scenarios}

The framework reproduces four CRS phenotypes via $\{\gamma_t\}$ and $\{K_f^{\text{scale}}_t\}$: \textbf{Type~1} (acute cardiac decompensation $\to$ AKI), \textbf{Type~2} (chronic HFpEF $\to$ progressive CKD via renal stress feedback), \textbf{Type~4} (primary CKD $\to$ secondary HFpEF via volume/RAAS-mediated remodeling), and \textbf{Combined} (simultaneous deterioration). The decomposition (Sec.~\ref{sec:decomposition}) quantifies each organ's decline attributable to its own disease versus cross-organ coupling.

{\small
\bibliographystyle{plainnat}
\bibliography{references}
}

\end{document}
